\section{Vocabulaire}

\begin{definitionbox}
\begin{itemize}[label=\textbullet]
    \item On appelle \bb{population statistique} l'ensemble 
        des \bb{individus} sur lesquels porte l'étude statistique. 
    \item On appelle \bb{effectif total} de la population statistique le
        nombre d'individus de cette population. 
    \item On appelle \bb{variable statistique} ou \bb{caractère}, la chose que l'on étudie
        et qui est commune à tous les individus de la population. L'ensemble des
        résultats s'appelle série statistique.
    \item On appelle \bb{effectif} associé à une valeur de la variable, le nombre de
        fois où cette valeur apparait dans la série.
\end{itemize}
\end{definitionbox}

\begin{exemplebox}
    Pour se rendre au collège

    \begin{itemize}[label=\textbullet]
        \item 71 élèves utilisent un deux roues,
        \item 305 élèves utilisent les transports en commune,
        \item 117 élèves se déplacent à pieds,
        \item 92 élèves sont déposés par leurs parents,
    \end{itemize}
    \begin{center}\rule{0.7\textwidth}{0.2pt}\end{center}
    On peut alors dire que :
    \begin{itemize}[label=\textbullet]
        \item la population statistique est \answer{les élèves}
        \item un individu est \answer{un élève}
        \item la variable (le caractère) est \answer{le moyen de transport utilisé}
        \item l'effectif des individus venant à pied est \answer{117}
    \end{itemize}

\end{exemplebox}

